\section{Results}
\subsection{Status-linked differences in the frequency of culture-talk}
Table~\ref{tab:talk-tab1} shows the point estimates and associated $t$-statistics corresponding to the predictors of expected self-reported frequency of culture talk across all five domains. First, consistent with hypothesis 1, it is clear that there exist strong levels of status-linked heterogeneity in culture-talk even after holding constant self-reported consumption behavior.  Consistent with the argument that culture talk is primarily stratified by socio-economic status, the results show that in all models education is among the best (if not the best) predictor of culture talk across all cultural domains. Persons who have attained advanced educational qualifications are more likely to engage in culture talk featuring popular culture and the arts as content than persons who have not, within levels of self-reported participation frequency. 


Second, for three of the five domains the effect of income is positive and statistically indistinguishable from the effects of education. This applies for the case of literature, the visual arts and drama ($p > 0.05$ for the likelihood ratio test constraining these coefficients to be the same). However, the effect of income is statistically smaller than the effect of education for movies ($p < 0.08$) and music ($p < 0.05$).  This suggests that persons who have access to economic resources are more likely to engage in culture talk than those who do not, but that education is a stronger differentiating factor in the popular and commercial realms. This is consistent with extant theory on the emerging cultural eclecticism of the educated elite \citep{peterson05, dimaggio96, ollivier04}, which points to the role of the commercial arts as a baseline conversational feature among educated persons complementing their long-standing command of traditional high-status cultures \citep{lizardo06a, dimaggio87}. In all, this pattern of results provides partial support for the hypothesis that the ability to translate consumption experiences into an interactional resource is more strongly dependent on the cultural resources obtained via experience in schooling institutions in comparison to economic resources \citep{bourdieu84, dimaggio96, lizardo06a}.

How does the additive model compare to the multiplicative specification?  Figure~\ref{fig:educgraph} shows the effect of education on the marginal probability of engaging in culture ``very often'' for respondents partitioned according to the levels of self-reported consumption in that domain.  The estimates are obtained from ordered probit models specified according to equation~\ref{eq:eq2}. Hypothesis 2 predicts that the education effect should be more pronounced among culturally active persons.  The results provide fairly strong levels of support for this hypothesis. Upward trends in the effect of education are clear in film book-reading, and to a lesser extent music.  They are much weaker in the case of dramatic performances. In the case of the visual arts, I simply find a categorical distinction between the most active participants and everybody else, with education-based differentiation in culture talk being strong for the former group and non-existent for the latter.  In all, these results provide fairly broad support for the hypothesis that status linked differences in culture talk are strong and are larger for the most active culture consumers.

Hypothesis 3 predicts that status-based differentiation should be stronger for the traditional high-status cultural domains, while hypothesis 4 predicts that we should find status-based differentiation in all domains.  A quick glance at the pattern of results shown in table~\ref{tab:talk-tab1} reveals that hypothesis 4 is on the right track.  The education advantage in culture talk is equally strong across all cultural domains.


  


\begin{comment}
%Third, older persons evince an invariably lower propensity of engaging in culture talk across all five domains, pointing to the importance of generational divisions in generating this outcome.  The age effects are strongest for culture talk involving film, and are weakest for the visual arts.  Put simply, older people are categorically less likely to use culture in interaction than younger people, net of consumption frequency and within levels of education and income.  This indicates that even when still sustaining an active participation profile, older persons tend to lose opportunities to translate those experiences into a resource in conversation.  One likely culprit is certainly the declining opportunities to interact with ``elective'' (non-kin) members of the personal network in settings conducive to culture talk that come with advanced age \citep{burt91}.  For instance, research on the personal networks of the elderly shows that personal networks decline in size and change in composition as persons age, with an increase in proportion kin and a decline in the proportion of non-kin contacts \citep{burt91, marsden87}.  Fourth, the positive coefficient estimate associated with gender is comparable in magnitude to the effect of income (for movies) and comparable to the effects of education (for literature, the visual arts and drama); there are no gender differences in culture talk propensity when it comes to music.  Note that this difference cannot be pegged to gender differences in cultural participation, since I am holding constant consumption frequency (and there are no systematic differences in the ways that men and women assign qualitative cutpoints to frequency categories as shown in table~\ref{tab:freq-tab2} in the Appendix).  Instead, it appears that women are systematically more likely to include aesthetic and literary experiences---for the cultural domains under consideration here---as a topic in conversation than men, suggesting that the literature, film and the fine arts are ``gendered'' via routine, interactional mechanisms \citep{ridgeway_smithlovin99}; research on the highly gendered nature of ``book clubs'' and their status as essentially a conversational vehicle \citep{long03}, as well as Erickson's \citeyearpar{erickson96} research noting the gendered status of the fine and performance arts in male-dominated workplaces \citep{lizardo06b}, is consistent with this hypothesis.\footnote{Presumably,  we can speculate that if we were investigating cultural domains tilted towards pursuits traditionally associated with masculinity (e.g. professional sports) we would find the opposite effect.}  Fifth, urban location is a good predictor of culture talk for four of the five domains under consideration here (with arts being the exception).  This suggests that opportunities to mobilize culture in interaction are facilitated by being embedded in contemporary urban contexts characterized by relative ease of access to leisure arenas where culture becomes a required type of social currency.


\subsection{Multiplicative versus additive specification}
Table~\ref{tab:talk-tab2} shows the coefficient estimates for the interaction between reported consumption frequency and selected predictors.  For the most part, I find that with the exception of film, the additive specification is to be preferred to the multiplicative one.  In the case of culture talk about the movies, I find that the high-education frequent consumers tend to translate this experience into higher rates of culture talk than their less educated counterparts.  When it comes to the age by frequency of consumption interaction for the same outcome, the results are best interpreted as negating the main---negative---effect of age for frequent-consumers.  However, the paucity of interaction effects in the other domains, and the virtue of parsimony in the face of very weak evidence \citep{raftery95}, strongly point to the superiority of the additive model.\footnote{Note that according to this criterion (BIC) only the multiplicative model in which culture talk involving film is the response variable evinces an improvement in fit over the additive model.}

In addition to including an interaction between frequency and socio-demographic predictors, I follow a suggestion inspired by recent research on SES-linked heterogeneity in the personal networks of older persons \citep{cornwell_etal08} and include an interaction between two key socio-demographic predictors: age and education.  Consistent with this research, a multiplicative model including this effect strongly recommends itself over the additive model that excludes this input for all five culture talk outcomes.  This suggests that cross-cutting cleavages premised on access to cultural resources exacerbate differences in the propensity to engage in culture talk across member of different status groups of the same generation.  As can be appreciated, education serves to negate a good chunk of the negative effect of age in the case of movies and music (where the negative age effect is greatest) and in the case of books, the visual arts and drama it serves to entirely negate the age effect.  This suggests that high-status older persons retain relatively larger propensities to engage in culture talk than their less-educated counterparts of the same age.  Another way of saying this is that the negative age-gradient in cultural talk propensities is much less negative (or even flat) among high-education groups than it is among low education groups; a finding that is consistent with observations of SES-related heterogeneity in network size among older adults  \citep{cornwell_etal08}. 

\subsection{Socio-demographic differences in cutpoint location}
Table~\ref{tab:talk-tab3} shows the coefficient estimates corresponding to the socio-demographic structuring of the qualitative location of the cutpoints on the underlying frequency distribution of culture talk. Recall that for those persons that perceive culture talk as normative for their peer group, we should find a right-ward shift in cutpoint location (indicated by a positive coefficient estimate), and the reverse for those groups who perceive culture talk as a less-pervasive feature of their social environment.  Following this line of reasoning, we should expect right-ward shift for instance, among educated, urban, wealthier elites.\footnote{To avoid over-parameterizing the cutpoint equations, I followed the following specification strategy:  I specified generalized ordered probit models predicting the (qualitatively indicated) frequency of consumption---the coefficient estimates corresponding to which are shown in Tables~\ref{tab:freq-tab1} and~\ref{tab:freq-tab2} in the Appendix---for each of the two cutpoint equations corresponding to each cultural domain, I only included those socio-demographic markers that proved to be significant predictors of consumption frequency for that domain.}  

The results are broadly consistent with this expectation.  Highly educated respondents consider higher underlying frequencies to be deserving of the category ``very often'' than low-education respondents for all five cultural domains.\footnote{Recall that the coefficient estimates are both centered and normalized, so that they retain an intuitive interpretation; thus in the case of movies, the expected cut point is 0.44 units (in the metric of the underlying frequency variable) higher for respondents in the highest education group in comparison to those in the lowest education group.} Education-related cutpoint shift for the category of ``fairly often'' is weaker, but can be found for both music and movies.  High-income respondents have tougher standards than low education respondents for considering a given frequency to be worthy of the ``fairly'' and ``very often'' labels in the case of movies, literature and the visual arts.  

Overall, there is very little residual heterogeneity in cutpoint location premised on age, gender and urban location after adjusting for education and income.  This suggests that the bulk of cross-respondent heterogeneity in the interpretation of qualitative frequency categories is driven by the most unambiguous markers of socio-economic status included in the model.  This is consistent with the the hypothesis that these are largely reflective of status group linked perceptions of the normative status of culture talk as a routine practice. 

\subsection{Summary}
This analysis reveals an intuitive picture of the ideal typical person for whom the mobilization of cultural knowledge in interaction is a routine activity:  Relatively urban, younger members of high-status elites (having ready access to high-levels of both cultural and economic resources).  These are the same groups who perceive culture talk as a relatively normative activity, as indicated by their propensity to consider underlying frequencies that their less-privileged counterparts deem to be ``very'' frequent as not quite deserving of that label.  Conversely, we can picture those who are least likely to engage in culture talk featuring aesthetic consumption experiences as a resource:  older, less-educated members of working-class or working poor status, cutting across ethnic and marital status designations.  The first group are almost a prototypical depiction of the ``new'' cosmopolitan elites that have been the subject of theorizing in the cultural capital and ``reading class'' literatures \citep{dimaggio_mukhtar04, katzgerro99, vaneijck00, dimaggio96, griswold08, emmison03}; the last are precisely those groups who suffer from the least access to prestigious symbolic and material goods, and which are thus excluded (and exclude themselves) from precisely those sites where elite cultural resources are accessed and disseminated \citep{bryson97, bourdieu_darbel91, blasius_friedrichs08, bennett_etal08}.  
\end{comment}